\documentclass[CJK]{beamer}
\usepackage{CJKutf8}
\usepackage{beamerthemesplit}
\usetheme{Malmoe}
\useoutertheme[footline=authortitle]{miniframes}
\usepackage{amsmath}
\usepackage{amssymb}
\usepackage{graphicx}
\usepackage{color}
\usepackage{slashed}
\usepackage{simplewick}
\graphicspath{{../figures/}}
\def\be{\begin{equation}}
\def\ee{\nonumber\end{equation}}
\def\bea{\begin{eqnarray}}
\def\eea{\nonumber\end{eqnarray}}
\def\ii{{\dot{\imath}}}
\def\bch{\begin{CJK}{UTF8}{gbsn}}
\def\ech{\end{CJK}}
\def\bex{\begin{minipage}{0.3\textwidth}\includegraphics[width=1in]{jugelizi.png}\end{minipage}\begin{minipage}{0.6\textwidth}}
\def\eex{\end{minipage}}
\def\chtitle#1{\frametitle{\bch#1\ech}}
\def\skipline{{\vskip0.1in}}
\def\skiplines{{\vskip0.2in}}
\def\lagr{{\mathcal{L}}}
\def\hamil{{\mathcal{H}}}
\def\vecv{{\mathbf{v}}}
\def\vecx{{\mathbf{x}}}
\def\veck{{\mathbf{k}}}
\def\vecp{{\mathbf{p}}}
\def\vecn{{\mathbf{n}}}
\def\vecA{{\mathbf{A}}}
\def\vecP{{\mathbf{P}}}
\def\vecsigma{{\mathbf{\sigma}}}
\def\hatJn{{\hat{J_\vecn}}}
\def\hatJx{{\hat{J_x}}}
\def\hatJy{{\hat{J_y}}}
\def\hatJz{{\hat{J_z}}}
\def\hatj#1{\hat{J_{#1}}}
\def\hatphi{{\hat{\phi}}}
\def\hatq{{\hat{q}}}
\def\hatpi{{\hat{\pi}}}
\def\vel{\upsilon}
\def\Dint{{\mathcal{D}}}
\def\adag{{\hat{a}^\dagger}}
\def\bdag{{\hat{b}^\dagger}}
\def\cdag{{\hat{c}^\dagger}}
\def\ddag{{\hat{d}^\dagger}}
\def\hata{{\hat{a}}}
\def\hatb{{\hat{b}}}
\def\hatc{{\hat{c}}}
\def\hatd{{\hat{d}}}
\def\hatN{{\hat{N}}}
\def\hatH{{\hat{H}}}
\def\hatp{{\hat{p}}}
\def\Fup{{F^{\mu\nu}}}
\def\Fdown{{F_{\mu\nu}}}
\def\newl{\nonumber \\}
\def\SIkm{\mathrm{km}}
\def\SIyr{\mathrm{yr}}
\def\SIGyr{\mathrm{Gyr}}
\def\SIeV{\mathrm{eV}}
\def\SIGeV{\mathrm{GeV}}
\def\SIm{\mathrm{m}}
\def\SIcm{\mathrm{cm}}
\def\SIJ{\mathrm{J}}
\def\SIs{\mathrm{s}}
\def\SIkg{\mathrm{kg}}
\def\SIg{\mathrm{g}}
\def\vece{\mathrm{e}}
\def\bmat#1{\left(\begin{array}{#1}}
\def\emat{\end{array}\right)}
\def\bcase#1{\left\{\begin{array}{#1}}
\def\ecase{\end{array}\right.}
\def\calM{{\mathcal{M}}}
\def\calT{{\mathcal{T}}}
\def\calR{{\mathcal{R}}}
\def\barpsi{\bar{\psi}}
\def\baru{\bar{u}}
\def\barv{\bar{\upsilon}}
\def\bmini#1{\begin{minipage}{#1\textwidth}}
\def\emini{\end{minipage}}
\def\qeq{\stackrel{?}{=}}
\def\torder#1{\mathcal{T}\left(#1\right)}
\def\rorder#1{\mathcal{R}\left(#1\right)}


\title{Quantum Field Theory I \\ Lesson 06 - Photon}
\author{}
\date{}


\begin{document}

\begin{frame}
 
\begin{center}
\begin{Large}
\bch
量子场论 I 

{\vskip 0.3in}

第七课 光子

\ech
\end{Large}
\end{center}

\vskip 0.2in

\bch
课件下载
\ech
https://github.com/zqhuang/SYSU\_QFTI

\end{frame}



\begin{frame}
\chtitle{$U(1)$规范场的量子化}
\bch
上节课我们留了一个问题还没有解决：自由$U(1)$规范场的粒子是怎样完成量子化的，这种粒子和谐振子又有何不同？

\skipline
这个问题的主要难点是每个时空点的$A^\mu$场仅有两个物理自由度，一旦我们要对这两个自由度进行量子化，就要取特定的规范来固定这两个自由度。不难想象，这样的推导并不会十分“优美”。

\ech
\end{frame}

\begin{frame}
\chtitle{第一种方案：库仑规范}
\bch
先试试库仑规范$A^0 = 0$, $\nabla \cdot \vecA = 0$.
这样$F_{00} = 0$, $F_{0i} = \partial_0 A_i$

$$ \lagr = -\frac{1}{4} \Fup\Fdown = \frac{1}{2} |\dot\vecA|^2  - \frac{1}{2}|\nabla \times \vecA|^2$$

对固定时间的拉氏量进行傅立叶变换后得到：
$$L = \int d^3\veck \left[\frac{1}{2} |\dot\vecA_\veck|^2  - \frac{1}{2}|\veck \times \vecA_\veck|^2\right]$$
我们对一个固定的$\veck$进行研究。

\ech
\end{frame}

\begin{frame}
\chtitle{第一种方案：库仑规范}
\bch
对固定的$\veck$，我们可以取三个方向$\vecn_1$, $\vecn_2$, $\vecn_3 = \veck / |\veck|$构成右手正交系。库仑规范要求$\vecA\cdot\veck = 0$, 所以$\vecA$可分解为$\vecA = u_{\veck,+1} \vece_{\veck,+1} + u_{\veck,-1} \vece_{\veck,-1}$。其中$\vece_{\veck,\pm 1} = \frac{1}{\sqrt{2}}(\vecn_1\pm \ii \vecn_2)$是$\hat{J}_{\vecn_3}$的本征值为$\pm 1$的归一化本征矢（参考上节课角动量的内容，添加了$1/\sqrt{2}$因子是为了归一化）。

利用正交归一条件$\vece_{\veck,+1}^\dagger \cdot \vece_{\veck,+1} = \vece_{\veck,-1}^\dagger \cdot \vece_{\veck,-1} = 1$, $\vece_{\veck,+1}^\dagger \cdot \vece_{\veck,-1} = \vece_{\veck,-1}^\dagger \cdot \vece_{\veck,+1} = 0$，可以得到：

$$|\dot\vecA_\veck|^2 = \dot\vecA_\veck^\dagger \cdot \dot\vecA_\veck = |\dot u_{\veck,+1}|^2 + |\dot u_{\veck,-1}|^2$$

\ech
\end{frame}


\begin{frame}
\chtitle{第一种方案：库仑规范}
\bch
注意到$\veck \times \frac{1}{\sqrt{2}}(\vecn_1\pm \ii\vecn_2) =|\veck|  \frac{1}{\sqrt{2}}(\vecn_2 \mp \ii \vecn_1)$, 而$ \frac{1}{\sqrt{2}}(\vecn_2 \mp \ii \vecn_1)$仍然是正交归一的两个矢量，于是得出
$$ |\veck \times \vecA_\veck|^2 = |\veck|^2(|u_{\veck,+1}|^2 + |u_{\veck,-1}|^2)$$

最终得到$\vecA_\veck$对拉氏量的贡献为:
$$ L_\veck = \frac{d^3\veck}{2}\sum_{s=\pm 1}  |\dot u_{\veck,s}|^2  - |\veck|^2 |u_{\veck,s}|^2$$

\ech
\end{frame}


\begin{frame}
\chtitle{第一种方案：库仑规范}
\bch

我们约定如果对$\veck$取了右手正交系$\vecn_1$, $\vecn_2$, $\veck/|\veck|$，则对$-\veck$就取$\vecn_1$, $-\vecn_2$, $-\veck/|\veck|$为正交右手系。那么对$s=\pm 1$都有
$$\vece_{-\veck, s}  = \vece_{\veck, s}^*$$

这样$\vecA$是实数场，也就是$\vecA_{-\veck} = \vecA_\veck^*$的条件就转化为$u_{-\veck, s} = u_{\veck, s}^*$。也就是说$u_{\veck, s}$ 可以看成一个实数场$u_s(\vecx)$的傅立叶变换。
\ech
\end{frame}

\begin{frame}
\chtitle{第一种方案：库仑规范}
\bch

这样，$u_s$的拉氏量就和质量为零的实标量场的拉氏量完全一样了。我们可以直接写出
$$ u_{\veck,s} = \frac{\hata_{\veck,s} + \adag_{-\veck,s}}{\sqrt{2 |\veck| d^3\veck}}$$。

最终进行反傅立叶变换，得到$\hat{\vecA}$的解为:
$$\hat{\vecA}(x) = \frac{1}{(2\pi)^{3/2}} \int \sqrt{\frac{d^3\veck}{2|\veck|}} \sum_{s=\pm 1}\mathbf{e}_{\veck s}\left(\hata_{\veck s} e^{-ik_\mu x^\mu} + \adag_{\veck s}e^{ik_\mu x^\mu}\right) $$
其中$\vece_{\veck,\pm 1} = \vecn_1 \pm \ii \vecn_2$, $\vecn_1$, $\vecn_2$, $\veck/|\veck|$是构成正交右手系的单位矢量，并如前所述有约定$\vece_{-\veck, s} = \vece_{\veck,s}^*$。 

\ech
\end{frame}

\begin{frame}
\chtitle{第一种方案：库仑规范}
\bch
上述推导中，我们顺便证明了U(1)规范场的粒子（例如光子）的质量为零，自旋为$1$。
\ech
\end{frame}

\begin{frame}
\chtitle{第二种方案：重矢量场}
拉氏密度里加质量项：
$$\lagr
\end{document}
